\documentclass[12pt, a4paper]{article}

\usepackage[T1]{fontenc}
\usepackage[utf8]{inputenc}
\usepackage[english]{babel}
\usepackage{geometry}
\usepackage{hyperref}
\usepackage{enumitem}
\usepackage{booktabs}
\usepackage{setspace}
\usepackage{parskip}
\usepackage{titlesec}
\usepackage[numbers,square,sort&compress]{natbib}

\geometry{margin=2.5cm}

\hypersetup{
    colorlinks=true,
    linkcolor=blue,
    citecolor=blue,
    urlcolor=blue
}

\title{\textbf{Design of Complex CI Pipelines}\\[0.5em]
\Large Registered Report}

\author{
    Brice Delcroix \and Timoth\'{e} Piscaglia \and Pierre-Louis Lef\`{e}vre\\[0.5em]
    \small Vint Cerf, INFOB302 -- IDS (2025--2026)\\
    \small Faculty of Computer Science, University of Namur, Belgium
}

\date{}

\begin{document}

\maketitle

% ──────────────────────────────────────────────────────────────
\section{Introduction}
% ──────────────────────────────────────────────────────────────

Continuous Integration (CI) has become a cornerstone of modern software engineering.
With platforms like GitHub Actions alone hosting millions of workflow definitions,
the scale and complexity of CI configurations have grown considerably.
By automatically building and testing code every time a developer commits changes
to a shared repository, CI ensures rapid feedback on code quality and drastically
reduces the cost of fixing defects~\cite{kinsman2021}.
As software projects grow in size and complexity, the CI pipelines that support
them inevitably grow as well. Today, pipelines may span multiple stages, multiple
language versions and operating systems, complex caching strategies, and integrations
with external quality-assurance services all encoded in declarative configuration
files~\cite{beller2017, dutta2025}.

Inefficiently designed or misconfigured pipelines can negate the very benefits that
CI is supposed to bring: slow feedback loops, frequent build failures, and high
maintenance burdens are well-documented consequences of poor pipeline design~\cite{dutta2025, hilton2016}.
Despite the widespread adoption of CI in both open-source and industrial contexts,
there is still no unified, cross-platform understanding of the configuration practices
and anti-patterns that characterise good versus poor CI pipeline design.
Existing empirical studies have identified anti-patterns for individual tools Travis
CI~\cite{gallaba2020}, GitLab CI~\cite{vassallo2020b}, GitHub Actions~\cite{kinsman2021} but
these findings remain specific to individual platforms and have not yet been consolidated
into a unified, cross-platform perspective on CI pipeline quality.

\paragraph{Scope:}
This research focuses on the design, architecture, and configuration management of
complex CI pipelines in open-source software projects, specifically on the three
most widely-used CI platforms: Travis CI, GitLab CI/CD, and GitHub Actions.
Using the PICOT framework to delimit the research:

\begin{itemize}
    \item \textbf{Population (P):} open-source software projects hosted on GitHub
          or GitLab that use Travis CI, GitLab CI, or GitHub Actions as their primary
          CI platform, and that have been continuously active (at least one pipeline
          execution per month) over the period 2018--2024.

    \item \textbf{Intervention (I):} systematic analysis of CI configuration files
          using existing anti-pattern catalogues and automated detection tools
          (Gretel~\cite{gallaba2020}, CI-Odor~\cite{vassallo2019}, CD-Linter~\cite{vassallo2020b}),
          combined with a manual labelling step to establish a ground truth.

    \item \textbf{Comparison (C):} comparison of anti-pattern prevalence rates across
          the three CI platforms and comparison of detection tool performance
          (precision, recall, coverage) against the manually labelled ground truth.

    \item \textbf{Outcome (O):} a cross-platform ranked catalogue of configuration
          anti-patterns and smells, a benchmarked assessment of each tool's detection
          capabilities, and an identification of quality issues not covered by any
          existing tool.

    \item \textbf{Time (T):} CI configuration files and repository histories from the
          period 2018--2024, covering the emergence and maturation of GitHub Actions
          (introduced in 2018), ensuring representation of all three platforms in
          their current form.
\end{itemize}

The research does not address continuous deployment (CD) aspects beyond CI scope,
nor human or organisational factors such as team adoption barriers.

\paragraph{Objectives:} 
This research pursues three specific objectives. First, it aims to identify and rank
the most prevalent configuration anti-patterns and smells across Travis CI, GitLab
CI, and GitHub Actions, extending platform-specific catalogues into a unified
cross-platform view. Second, it aims to evaluate how well existing automated
detection tools (Gretel, CI-Odor, CD-Linter) cover the identified anti-patterns,
measuring their precision, recall, and coverage gaps. Third, it aims to identify
categories of configuration quality issues that remain unaddressed by current tooling,
thereby guiding future tool development. In terms of knowledge gain, the expected
empirical contribution is a benchmarked comparison of tool effectiveness on a
manually labelled dataset; the expected theoretical contribution is a cross-platform
anti-pattern taxonomy that consolidates and extends prior single-platform catalogues.
These contributions are of practical relevance to CI tool developers and DevOps
engineers seeking to improve pipeline quality assessment, and of societal relevance
in so far as more reliable CI pipelines directly translate to more reliable software
delivered to end-users.

\paragraph{State-of-the-art positioning:}
As detailed in Section~3, the existing literature has produced single-platform
anti-pattern catalogues~\cite{gallaba2020, vassallo2020b, kinsman2021} and individual
tool evaluations~\cite{gallaba2020, vassallo2019, vassallo2020b}, but no study has
evaluated the three major detection tools side-by-side against a common benchmark.
This research fills that gap by adopting a cross-platform, tool-comparative methodology.

\paragraph{Research method overview:}
We follow an empirical research strategy combining large-scale repository mining (to
establish a cross-platform anti-pattern inventory) with a controlled benchmarking
evaluation (to assess tool performance). The mining protocol draws on the methodology
of Gallaba and McIntosh~\cite{gallaba2020} and Vassallo et al.~\cite{vassallo2020b},
while the benchmarking approach follows principles from controlled empirical studies
in software engineering~\cite{kinsman2021}.

% ──────────────────────────────────────────────────────────────
\section{Background}
% ──────────────────────────────────────────────────────────────

This section introduces the key concepts required to understand the scope and
methodology of this research.

\subsection{Continuous Integration Pipelines}

Continuous Integration is the practice of automatically building and testing software
each time a developer pushes changes to a version control repository~\cite{kinsman2021}.
A CI pipeline is the set of automated steps that implement this practice, typically
including dependency management, compilation, testing, static analysis, and artefact
generation. Modern CI tools Travis CI, GitLab CI/CD, GitHub Actions, and CircleCI encode
these steps in declarative configuration files (e.g., \texttt{.travis.yml},
\texttt{.gitlab-ci.yml}, or GitHub Actions YAML workflows)~\cite{bozhkov2024, gallaba2020}.
As projects scale, these files grow in complexity: they may involve matrix builds
across multiple language versions and operating systems, conditional job execution,
caching strategies, and integrations with external services. St{\aa}hl and Bosch~\cite{stahl2014}
proposed a descriptive model capturing the key variation points in CI practices integration
scope, trigger mechanisms, and build frequency providing a useful vocabulary for
reasoning about pipeline design diversity across industrial contexts.

\subsection{Configuration as Code}

Configuration as Code (CaC) is the practice of managing CI pipeline definitions in
version-controlled text files, with the same rigour applied to production source
code~\cite{layman2022}. This approach enables reproducibility, auditability, and
collaborative review of pipeline configurations and, crucially, it makes
configurations amenable to static analysis by automated tools. Nonetheless, Nejati
et al.~\cite{nejati2023} showed that build system specifications are reviewed at
roughly half the frequency of production code, yet review comments are more likely
to identify defects. This suggests that pipeline configurations are under-scrutinised
despite their critical role, reinforcing the need for automated quality assessment.

\subsection{Configuration Anti-Patterns and Smells}

By analogy with code smells in traditional software engineering, configuration smells
are suboptimal patterns in pipeline configuration files that may indicate deeper
design problems~\cite{gallaba2020}. Gallaba and McIntosh~\cite{gallaba2020} identified
four categories of anti-patterns in Travis CI specifications, finding that nearly half
of all configuration code is devoted to environment setup rather than core
build-and-test logic. Vassallo et al.~\cite{vassallo2020b} defined four types of
configuration smells in GitLab CD pipelines and implemented CD-Linter, achieving 87\%
precision and 94\% recall over 5,312 projects. Zampetti et al.~\cite{zampetti2020}
compiled the most exhaustive taxonomy to date 79 bad practices grouped into seven
categories based on developer interviews, Stack Overflow analysis, and surveys.
These studies collectively demonstrate that configuration quality issues are pervasive,
diverse, and platform-dependent.

\subsection{Automated Detection Tools}

Several tools have been proposed to detect quality issues in CI configurations
automatically. Gretel, introduced by Gallaba and McIntosh~\cite{gallaba2020}, detects
anti-patterns in Travis CI configuration files. CI-Odor, proposed by Vassallo et
al.~\cite{vassallo2019}, operates at a different level of abstraction: rather than
analysing configuration files directly, it identifies four process-level anti-patterns
(Slow Build, Broken Release, Late Merging, and Missing Build Failure Repair) by mining
build logs and repository metadata. This distinction is important for our benchmarking
protocol, as CI-Odor's detection targets differ in nature from those of Gretel and
CD-Linter. CD-Linter~\cite{vassallo2020b} targets configuration smells in GitLab
pipeline YAML files. While each tool has shown promising results on its target
platform, it remains unknown whether their detection capabilities generalise across
CI systems and whether important categories of design issues remain
unaddressed---the central gap this research addresses.

\subsection{Build Failures and Pipeline Quality Consequences}

A body of empirical research documents the consequences of poor pipeline design.
Rausch et al.~\cite{rausch2017} analysed 4,021 Travis CI builds across 30 Java
projects, categorising failure causes into compilation errors, test failures, and
environment issues. Beller et al.~\cite{beller2017} showed that most build breaks
are caused by test failures rather than compilation errors, and that overnight
failures take significantly longer to fix. Vassallo et al.~\cite{vassallo2017}
compared build failure patterns across open-source and financial-sector contexts,
revealing domain-specific failure categories. These findings motivate the need for
early, automated detection of configuration quality issues.

% ──────────────────────────────────────────────────────────────
\section{Related Work}
% ──────────────────────────────────────────────────────────────

Our research sits at the intersection of empirical studies on CI configuration quality
and tool-based research on automated pipeline analysis. We position our work explicitly
with respect to the most relevant related research.

\subsection{Same Problem, Complementary Approaches (Single-Platform Studies)}

Several studies have investigated configuration quality issues in CI pipelines, but
each has focused on a single platform. Gallaba and McIntosh~\cite{gallaba2020} studied
anti-patterns exclusively in Travis CI, identifying four anti-pattern categories in
9,312 projects and proposing the Gretel detection tool. Vassallo et al.~\cite{vassallo2020b}
targeted GitLab CI/CD, defining four configuration smell types and implementing
CD-Linter. Kinsman et al.~\cite{kinsman2021} provided the first empirical
characterisation of GitHub Actions workflows but did not specifically analyse
configuration quality. Our research differs from all three by adopting a cross-platform
perspective, comparing anti-pattern prevalence across Travis CI, GitLab CI, and GitHub
Actions to identify platform-specific versus universal design issues. This cross-platform
synthesis is the primary novelty of our contribution relative to this body of work.

Table~\ref{tab:related} summarises how the most directly related studies compare to
our work along the dimensions that define our novelty.

\begin{table}[h]
\centering
\caption{Comparison of related studies on CI configuration quality.}
\label{tab:related}
\footnotesize
\begin{tabular}{p{3.2cm}p{2.8cm}p{2.4cm}p{3.6cm}p{1.4cm}}
\toprule
\textbf{Study} & \textbf{Platform(s)} & \textbf{Anti-patterns} & \textbf{Tool evaluated} & \textbf{Cross-platform?} \\
\midrule
Gallaba \& McIntosh~\cite{gallaba2020}   & Travis CI         & 4 categories    & Gretel                    & No      \\
Vassallo et al.~\cite{vassallo2020b}     & GitLab CI         & 4 smell types   & CD-Linter                 & No      \\
Vassallo et al.~\cite{vassallo2019}      & Platform-agnostic & 4 process-level & CI-Odor                   & Partial \\
Kinsman et al.~\cite{kinsman2021}        & GitHub Actions    & General workflows & None                    & No      \\
Zampetti et al.~\cite{zampetti2020}      & Multiple (survey) & 79 practices    & None                      & Partial \\
\textbf{This work}                       & \textbf{Travis, GitLab, GHA} & \textbf{79+ practices} & \textbf{Gretel, CI-Odor, CD-Linter} & \textbf{Yes} \\
\bottomrule
\end{tabular}
\end{table}

\subsection{Same Problem, Same Approach, Different Goal (Evolutionary Studies)}

Zampetti et al.~\cite{zampetti2021} conducted both qualitative and quantitative
analyses of pipeline evolution, producing a taxonomy of 34 restructuring operations
across 4,644 Travis CI projects. Their goal was to understand how pipelines change
over time, whereas our goal is to evaluate how well existing tools detect the quality
issues that motivate such changes. Similarly, Rausch et al.~\cite{rausch2017} and
Vassallo et al.~\cite{vassallo2017} studied build failure causes in Java projects
and across open-source versus financial-sector contexts, respectively. These studies
provide empirical evidence of the consequences of poor pipeline design but do not
assess detection tooling which is our central concern. The comprehensive bad-practice
taxonomy of Zampetti et al.~\cite{zampetti2020} serves as a reference catalogue for
our gap analysis, rather than as a direct comparative object.

\subsection{Similar Approach, Broader Problem (Systematic Reviews)}

Shahin et al.~\cite{shahin2017a} and Laukkanen et al.~\cite{laukkanen2017} produced
systematic literature reviews covering the full CI/CD landscape, including challenges,
tools, and practices. Our research is narrower in scope focusing specifically on
configuration-level quality but deeper in its evaluation of automated detection
tools. Shahin et al.~\cite{shahin2019} further examined architectural decisions
supporting CD pipelines through practitioner interviews, complementing our
configuration-centric analysis. Saleh et al.~\cite{saleh2024} reviewed CI/CD for
secure cloud computing, providing additional architectural context. These reviews
inform our background section but are not direct comparators for our empirical
contribution.

\subsection{Operational and Evolutionary Perspectives}

Gallaba et al.~\cite{gallaba2022} analysed eight years of CircleCI operational data
to study scalability trends, revealing how pipeline design decisions accumulate into
long-term performance consequences. Widder et al.~\cite{widder2018} investigated why
projects abandon Travis CI, identifying build duration and programming language as
significant predictors motivating the need for portable, tool-agnostic pipeline
quality assessment. These operational studies contextualise the long-term consequences
of the configuration quality issues our research aims to detect early. Zampetti et
al.~\cite{zampetti2017b} examined how open-source projects integrate static analysis
tools into CI pipelines, directly informing our understanding of quality-gate
governance in complex pipelines.

% ──────────────────────────────────────────────────────────────
\section{Research Question(s)}
% ──────────────────────────────────────────────────────────────

Based on the gaps identified in the state of the art, we formulate the following
main research question:

\bigskip
\noindent\textbf{RQ\quad} \textit{To what extent do existing automated tools detect
configuration anti-patterns and smells in complex CI pipelines across different CI
platforms?}
\bigskip

This question is motivated by the observation that current detection tools
(Gretel~\cite{gallaba2020}, CI-Odor~\cite{vassallo2019}, CD-Linter~\cite{vassallo2020b})
have each been designed and evaluated for a single CI platform, targeting a subset
of known anti-patterns. Whether their coverage is representative of the full spectrum
of configuration quality issues, and whether their detection capabilities transfer to
other platforms, remains unknown. Answering this question contributes both a
comparative evaluation of tool effectiveness and an identification of under-addressed
quality concerns. The question satisfies the FINER criteria: it is \emph{Feasible}
(relies on publicly available open-source repositories and existing tools),
\emph{Interesting} (no prior cross-platform comparison exists), \emph{Novel} (extends
single-platform catalogues), \emph{Ethical} (uses only public data), and
\emph{Relevant} (directly informs CI tool development and DevOps practice).

We decompose the main research question into three sub-research questions:

\paragraph{RQ.1}
\textit{What configuration anti-patterns and smells are most prevalent in CI pipelines
across Travis CI, GitLab CI, and GitHub Actions, and do their frequency and severity
differ by platform?}

This sub-question targets open-source projects active between 2018 and 2024 across
Travis CI, GitLab CI, and GitHub Actions. Through systematic mining and manual
labelling based on established anti-pattern definitions~\cite{gallaba2020, vassallo2020b, zampetti2020},
we aim to produce a ranked catalogue of anti-patterns by frequency and severity,
stratified by platform. This extends the single-platform catalogues of Gallaba and
McIntosh~\cite{gallaba2020}, Vassallo et al.~\cite{vassallo2020b}, and Kinsman et
al.~\cite{kinsman2021} by enabling direct cross-platform comparisons.

\paragraph{RQ.2}
\textit{How do existing automated detection tools (Gretel, CI-Odor, CD-Linter)
compare in terms of precision, recall, and coverage of the identified anti-patterns?}

This sub-question evaluates the detection capabilities of Gretel, CI-Odor, and
CD-Linter by applying each tool to the manually labelled ground-truth dataset
constructed for RQ.1 and measuring precision, recall, and F1-score per tool and per
anti-pattern category. While each tool has been individually validated on its target
platform~\cite{gallaba2020, vassallo2019, vassallo2020b}, no study has compared them
side-by-side on a common benchmark which is the central contribution of this
sub-question.

\paragraph{RQ.3}
\textit{What categories of CI pipeline configuration quality issues are not addressed
by any existing automated detection tool?}

This sub-question identifies the gaps left by current tooling by comparing the
cross-platform inventory from RQ.1 with the coverage results from RQ.2. Using the
taxonomy of Zampetti et al.~\cite{zampetti2020} (79 bad practices in 7 categories)
and the restructuring operations catalogued by Zampetti et al.~\cite{zampetti2021}
as references, we pinpoint which categories of configuration problems currently escape
automated detection. The expected outcome is a prioritised list of unaddressed quality
concerns to guide future tool development.

% ──────────────────────────────────────────────────────────────
\section{Evaluation Protocol}
% ──────────────────────────────────────────────────────────────

This section describes the evaluation protocol in sufficient detail for an independent
researcher to replicate it without making additional design decisions.

\subsection{Dataset Construction (for RQ.1 and RQ.2)}

\subsubsection{Repository Selection}

We will collect open-source repositories from GitHub and GitLab using the GitHub REST
API and the GitLab API. Repositories must satisfy all of the following inclusion
criteria:

\begin{itemize}
    \item The repository uses Travis CI (\texttt{.travis.yml} present), GitLab CI
          (\texttt{.gitlab-ci.yml} present), or GitHub Actions
          (\texttt{.github/workflows/} directory present) as its primary CI platform.

    \item The repository has been continuously active defined as at least one CI
          pipeline execution per month from January 2018 to December 2024.

    \item The repository has at least 100 stars on GitHub (or 50 on GitLab), as a
          proxy for project maturity and non-triviality.

    \item The repository is not a fork of another included repository.

    \item The primary language is one of: Java, Python, JavaScript, Ruby, or Go (to
          control for language-specific build patterns).
\end{itemize}

We target a stratified sample of 50 repositories per platform (150 total), stratified
by primary language (10 per language per platform). This sample size balances the
feasibility constraints of manual labelling with the need for statistical
representativeness. If fewer than 40 repositories satisfy all criteria for a given
stratum, we include all qualifying repositories. The sampling will use the GHTorrent
dataset~\cite{gousios2013} as a starting point, supplemented by direct API queries
for repositories created after the GHTorrent snapshot.

To validate the automated filtering, two team members will independently inspect a
random sample of 30 repositories drawn from the filtered results and confirm that each
genuinely satisfies all inclusion criteria. Any repository for which the two reviewers
disagree will be excluded by default. This manual check also serves to detect edge
cases in which the presence of a CI configuration file does not reflect active use
of that platform (e.g., disabled or empty workflows).

\subsubsection{Data Extraction}

For each repository, we will extract: (1) all CI configuration files, (2) the full
Git history of those configuration files, and (3) basic project metadata (number of
contributors, language, age, number of CI runs). Configuration files will be extracted
at three time points: January 2019, January 2022, and December 2024, capturing early,
middle, and current pipeline designs.

Before analysis, all extracted configuration files will undergo the following
preprocessing steps. First, files that cannot be parsed as valid YAML will be
flagged and excluded from the benchmarking sub-sample, though they will be counted
in the dataset to avoid inflating quality metrics. Second, duplicate files ---
defined as files with identical content appearing at multiple time points for the
same repository --- will be deduplicated, retaining only the most recent version
for the labelling step to avoid inflated anti-pattern counts. Third, all files will
be normalised to UTF-8 encoding and Unix line endings prior to tool execution, to
ensure consistent behaviour across the three detection tools. These steps will be
fully scripted and included in the replication package (Section~5.5).

\subsubsection{Manual Labelling (Ground Truth)}

To evaluate tool performance (RQ.2), we will construct a ground-truth dataset by
manually labelling a stratified sub-sample of 60 configuration files (20 per platform).
This sample size is consistent with prior manual annotation studies in CI research:
Vassallo et al.~\cite{vassallo2020b} manually validated a comparable set of files
when evaluating CD-Linter, and Gallaba and McIntosh~\cite{gallaba2020} used similar
manual inspection to validate Gretel's outputs. Given the exploratory nature of this
benchmarking study, which prioritises cross-platform breadth over statistical depth,
we acknowledge the limited size as a boundary condition and report effect sizes
alongside all statistical tests to compensate for reduced power.

Three team members will independently label each file using the anti-pattern catalogue
of Gallaba and McIntosh~\cite{gallaba2020}, Vassallo et al.~\cite{vassallo2020b},
and Zampetti et al.~\cite{zampetti2020} as reference. We acknowledge that the
labellers are also the study designers, which introduces a potential confirmation bias.
We mitigate this through three complementary measures: (1) blind labelling, where
each annotator works independently without knowledge of others' judgements; (2) the
use of a structured binary checklist derived directly from the published anti-pattern
definitions, reducing interpretive freedom to yes/no decisions per anti-pattern per
line; and (3) an audit of a randomly drawn subset of ten files by a fourth reviewer
not involved in the study design, whose agreement with the majority label will be
reported as an additional reliability indicator.
Each label records: anti-pattern type, location (line number), and severity
(minor/moderate/severe). Inter-rater agreement will be measured using Cohen's Kappa;
disagreements will be resolved by majority vote after discussion. We consider
$\kappa \geq 0.70$ as acceptable agreement.

\subsection{Anti-Pattern Inventory (RQ.1)}

On the full 150-repository dataset, we will apply existing anti-pattern definitions
from~\cite{gallaba2020, vassallo2020b, zampetti2020} to automatically identify
occurrences of known patterns where tool support exists, and manually annotate the
sub-sample for patterns without automated support. For each anti-pattern, we will
compute: (1) prevalence rate (percentage of repositories exhibiting it), (2) frequency
per file (mean occurrences per configuration file), and (3) platform distribution. We
will report descriptive statistics (median, IQR, minimum, maximum) and use
Mann--Whitney U tests (significance level $\alpha = 0.05$, Bonferroni-corrected for
multiple comparisons) to assess whether prevalence rates differ significantly between
platforms. Prior to applying these tests, we will verify that the two required
assumptions hold: that the data are at least ordinal, and that observations are
independent across platforms (satisfied by design, since each repository belongs to
exactly one platform stratum).

\subsection{Tool Benchmarking (RQ.2)}

We will apply Gretel and CD-Linter to the ground-truth sub-sample for
configuration-level detection. Since CI-Odor operates at the process level (mining
build logs and repository metadata rather than analysing configuration files directly),
it is fundamentally incommensurable with Gretel and CD-Linter on precision and recall
metrics defined over configuration-file annotations. CI-Odor will therefore be
evaluated separately using build history data from the same repositories, and its
results presented in a dedicated subsection with qualitative discussion; they will
\emph{not} be merged into the precision/recall/F1 comparison table used for Gretel
and CD-Linter. For Gretel and CD-Linter, we will measure:

\begin{itemize}
    \item \textbf{Precision:} proportion of tool-flagged instances that correspond to
          a true anti-pattern in the ground truth.

    \item \textbf{Recall:} proportion of ground-truth anti-pattern instances that are
          flagged by the tool.

    \item \textbf{F1-score:} harmonic mean of precision and recall.

    \item \textbf{Coverage:} proportion of anti-pattern categories (from the Zampetti
          et al.~\cite{zampetti2020} taxonomy) for which the tool flags at least one
          instance.
\end{itemize}

Results will be reported both aggregated (overall precision/recall/F1) and per
anti-pattern category. Note that Gretel targets Travis CI exclusively, CD-Linter
targets GitLab CI exclusively, and CI-Odor is platform-agnostic; we will therefore
apply each tool only to the configurations of its target platform when measuring
precision and recall, but measure coverage across all platforms.

\subsection{Gap Analysis (RQ.3)}

We will construct a coverage matrix mapping each of the 79 bad practices from Zampetti
et al.~\cite{zampetti2020} and the 34 restructuring operations from Zampetti et
al.~\cite{zampetti2021} to the detection tools. A category is considered ``covered''
if at least one of the three tools achieves recall $\geq 0.50$ for instances of that
category in the ground truth. This threshold follows the convention used by Vassallo
et al.~\cite{vassallo2020b}, who considered detection rates above 50\% as indicative
of meaningful coverage. Categories below this threshold, or absent from tool output
entirely, are classified as ``gaps''. We will prioritise the gap list by: (1)
prevalence in the RQ.1 dataset, (2) severity ratings from the manual labelling, and
(3) theoretical impact based on the literature (e.g., build failure
likelihood~\cite{rausch2017, vassallo2017}).

\subsection{Replication Package}

All scripts, sampled repository lists, manual labelling guidelines, labelled
ground-truth data, and raw tool outputs will be made publicly available on a dedicated
Zenodo archive upon completion of the study, to ensure full replicability.

\subsection{Threats to Validity}

\paragraph{Internal Validity.}
Subjectivity in manual labelling is the primary internal threat. We mitigate this by
using three independent raters and measuring inter-rater agreement with Cohen's Kappa,
resolving disagreements through discussion. To further reduce confirmation bias
arising from labellers also being study designers, we apply a structured binary
checklist derived from published anti-pattern definitions~\cite{gallaba2020, vassallo2020b, zampetti2020}
that limits interpretive freedom, and we include an audit by a fourth reviewer
external to the study design (see Section~5.1.3). Another threat is the possible
misapplication of a tool outside its intended platform; we mitigate this by applying
each tool only to its target platform for performance measurement (Section~5.3).

\paragraph{Construct Validity.}
We operationalise ``anti-pattern prevalence'' as the percentage of repositories
exhibiting an instance, which may not capture severity. We address this by also
reporting frequency per file and severity ratings. Precision and recall measure tool
performance against a human-labelled ground truth, which itself may be incomplete; we
mitigate this by using three raters and published taxonomies. The star threshold (100
stars) as a proxy for project maturity may exclude relevant smaller projects; we
acknowledge this as a boundary condition of our sampling strategy.

\paragraph{External Validity.}
The sample is limited to five programming languages and three CI platforms, and to
open-source repositories. Results may not generalise to proprietary CI systems or to
languages not included. The restriction to 2018--2024 may also limit applicability
to earlier or future pipeline configurations. We partially mitigate these threats by
including the three most widely-used CI platforms~\cite{gallaba2020, kinsman2021, vassallo2020b}
and by stratifying across languages. We explicitly report the scope constraints and
encourage replication in industrial settings.

\paragraph{Reliability.}
All tools (Gretel, CI-Odor, CD-Linter) are open-source and their versions will be
fixed and archived in the replication package. API queries for repository selection
will be scripted and logged. The full replication package (Section~5.5) enables
independent verification of all steps.

\paragraph{Conclusion Validity.}
With a ground-truth dataset of 60 configuration files (20 per platform), the
statistical power of pairwise comparisons may be limited. We mitigate this by
reporting effect sizes alongside $p$-values and by using non-parametric tests
appropriate for small samples. We acknowledge that Bonferroni correction increases
the risk of Type II errors at this sample size.

% ──────────────────────────────────────────────────────────────
\begin{thebibliography}{99}
% ──────────────────────────────────────────────────────────────

\bibitem{beller2017}
M.~Beller, G.~Gousios, and A.~Zaidman.
``Oops, my tests broke the build: An explorative analysis of Travis CI with GitHub.''
In \emph{Proc. 14th IEEE/ACM Int. Conf. on Mining Software Repositories (MSR)},
pp.~356--367. IEEE, 2017.

\bibitem{bozhkov2024}
D.~Bozhkov et al.
``Configuration tool for CI/CD pipelines and React web apps.''
In \emph{2024 14th Int. Conf. on Advanced Computer Information Technologies (ACIT)}.
IEEE, 2024.

\bibitem{cassee2020}
N.~Cassee, B.~Vasilescu, and A.~Serebrenik.
``The silent helper: The impact of continuous integration on code reviews.''
In \emph{Proc. 27th IEEE Int. Conf. on Software Analysis, Evolution and Reengineering
(SANER)}, pp.~423--434. IEEE, 2020.

\bibitem{dutta2025}
S.~Dutta et al.
``CI/CD pipeline optimization using AI: A systematic mapping study.''
\emph{MDPI Applied Sciences}, 2025.

\bibitem{fitzgerald2017}
B.~Fitzgerald and K.-J. Stol.
``Continuous software engineering: A roadmap and agenda.''
\emph{Journal of Systems and Software}, 123:176--189, 2017.

\bibitem{gallaba2022}
K.~Gallaba, M.~Lamothe, and S.~McIntosh.
``Lessons from eight years of operational data from a continuous integration service:
An exploratory case study of CircleCI.''
In \emph{Proc. 44th IEEE/ACM Int. Conf. on Software Engineering (ICSE)},
pp.~1330--1342. ACM, 2022.

\bibitem{gallaba2020}
K.~Gallaba and S.~McIntosh.
``Use and misuse of continuous integration features: An empirical study of projects
that (mis)use Travis CI.''
\emph{IEEE Trans. on Software Engineering}, 46(1):33--50, 2020.

\bibitem{hilton2016}
M.~Hilton, T.~Tunnell, K.~Huang, D.~Marinov, and D.~Dig.
``Usage, costs, and benefits of continuous integration in open-source projects.''
In \emph{Proc. 31st IEEE/ACM Int. Conf. on Automated Software Engineering (ASE)},
pp.~426--437. ACM, 2016.

\bibitem{karlovs2025}
M.~Karlovs-Karlovskis and O.~Nikiforova.
``From software architecture models to pipelines: A conceptual framework for model
transformation in DevOps.''
\emph{Frontiers in Computer Science}, 2025.

\bibitem{kinsman2021}
T.~Kinsman, M.~Wessel, M.A.~Gerosa, and C.~Treude.
``How do software developers use GitHub Actions to automate their workflows?''
In \emph{Proc. 18th IEEE/ACM Int. Conf. on Mining Software Repositories (MSR)},
pp.~420--431. IEEE, 2021.

\bibitem{laukkanen2017}
M.~Laukkanen, J.~Itkonen, and C.~Lassenius.
``Problems, causes and solutions when adopting continuous delivery: A systematic
literature review.''
\emph{Information and Software Technology}, 82:55--79, 2017.

\bibitem{layman2022}
L.~Layman et al.
``Method for continuous integration and deployment using a pipeline generator for
agile software projects.''
\emph{Sensors (MDPI)}, 22(12):4637, 2022.

\bibitem{leppanen2015}
M.~Lepp{\"a}nen et al.
``The highways and country roads to continuous deployment.''
\emph{IEEE Software}, 32(2):64--72, 2015.

\bibitem{nejati2023}
M.~Nejati, M.~Alfadel, and S.~McIntosh.
``Code review of build system specifications: Prevalence, purposes, patterns, and
perceptions.''
In \emph{Proc. 45th IEEE/ACM Int. Conf. on Software Engineering (ICSE)},
pp.~1213--1224. IEEE, 2023.

\bibitem{rausch2017}
T.~Rausch, W.~Hummer, P.~Leitner, and S.~Schulte.
``An empirical analysis of build failures in the continuous integration workflows of
Java-based open-source software.''
In \emph{Proc. 14th Int. Conf. on Mining Software Repositories (MSR)},
pp.~345--355. IEEE, 2017.

\bibitem{saleh2024}
S.M.~Saleh, N.H.~Madhavji, and J.~Steinbacher.
``A systematic literature review on continuous integration and deployment (CI/CD) for
secure cloud computing.''
In \emph{Proc. 20th Int. Conf. on Web Information Systems and Technologies (WEBIST)}.
SciTePress, 2024.

\bibitem{shahin2017a}
M.~Shahin, M.A.~Babar, and L.~Zhu.
``Continuous integration, delivery and deployment: A systematic review on approaches,
tools, challenges and practices.''
\emph{IEEE Access}, 5:1920--1938, 2017.

\bibitem{shahin2019}
M.~Shahin, M.~Zahedi, M.A.~Babar, and L.~Zhu.
``An empirical study of architecting for continuous delivery and deployment.''
\emph{Empirical Software Engineering}, 24(3):1061--1108, 2019.

\bibitem{stahl2014}
D.~St{\aa}hl and J.~Bosch.
``Modeling continuous integration practice differences in industry software
development.''
\emph{Journal of Systems and Software}, 87:48--59, 2014.

\bibitem{ullah2025}
A.~Ullah et al.
``Systematic mapping study on AI integration in CI/CD pipelines.''
2025. Preprint---peer-reviewed status to be confirmed.

\bibitem{vassallo2019}
C.~Vassallo, S.~Proksch, H.C.~Gall, and M.~Di~Penta.
``Automated reporting of anti-patterns and decay in continuous integration.''
In \emph{Proc. 41st IEEE/ACM Int. Conf. on Software Engineering (ICSE)},
pp.~105--115. IEEE, 2019.

\bibitem{vassallo2020b}
C.~Vassallo, S.~Proksch, A.~Jancso, H.C.~Gall, and M.~Di~Penta.
``Configuration smells in continuous delivery pipelines: A linter and a six-month
study on GitLab.''
In \emph{Proc. 28th ACM Joint European Software Engineering Conf. and Symposium on
the Foundations of Software Engineering (ESEC/FSE)}, pp.~327--337. ACM, 2020.

\bibitem{vassallo2017}
C.~Vassallo et al.
``A tale of CI build failures: An open source and financial sector analysis.''
In \emph{Proc. 33rd IEEE Int. Conf. on Software Maintenance and Evolution (ICSME)},
pp.~183--193. IEEE, 2017.

\bibitem{widder2018}
D.~Widder, B.~Vasilescu, M.~Hilton, and C.~K{\"a}stner.
``I'm leaving you, Travis: A continuous integration breakup story.''
In \emph{Proc. 15th Int. Conf. on Mining Software Repositories (MSR)},
pp.~165--169. IEEE, 2018.

\bibitem{wimmer2021}
M.~Wimmer et al.
``Continuous integration in multi-view modeling: A model transformation pipeline
architecture for production systems engineering.''
In \emph{Proc. 9th Int. Conf. on Model-Driven Engineering and Software Development
(MODELSWARD)}. SciTePress, 2021.

\bibitem{zampetti2019}
F.~Zampetti et al.
``Detecting architectural issues during the continuous integration pipeline.''
In \emph{2019 ACM/IEEE 22nd Int. Conf. on Model Driven Engineering Languages and
Systems Companion (MODELS-C)}. IEEE, 2019.

\bibitem{zampetti2021}
F.~Zampetti, S.~Geremia, G.~Bavota, and M.~Di~Penta.
``CI/CD pipelines evolution and restructuring: A qualitative and quantitative study.''
In \emph{Proc. 37th IEEE Int. Conf. on Software Maintenance and Evolution (ICSME)},
pp.~471--482. IEEE, 2021.

\bibitem{zampetti2017b}
F.~Zampetti, S.~Scalabrino, R.~Oliveto, G.~Canfora, and M.~Di~Penta.
``How open source projects use static code analysis tools in continuous integration
pipelines.''
In \emph{Proc. 14th Int. Conf. on Mining Software Repositories (MSR)},
pp.~334--344. IEEE, 2017.

\bibitem{zampetti2020}
F.~Zampetti, C.~Vassallo, S.~Panichella, G.~Canfora, H.C.~Gall, and M.~Di~Penta.
``An empirical characterization of bad practices in continuous integration.''
\emph{Empirical Software Engineering}, 25(2):1095--1135, 2020.

\bibitem{zhao2017}
Y.~Zhao, A.~Serebrenik, Y.~Zhou, V.~Filkov, and B.~Vasilescu.
``The impact of continuous integration on other software development practices:
A large-scale empirical study.''
In \emph{Proc. 32nd IEEE/ACM Int. Conf. on Automated Software Engineering (ASE)},
pp.~60--71. IEEE, 2017.

\bibitem{gousios2013}
G.~Gousios.
``The GHTorrent dataset and tool suite.''
In \emph{Proc. 10th Working Conference on Mining Software Repositories (MSR)},
pp.~233--236. IEEE, 2013.

\end{thebibliography}

% ──────────────────────────────────────────────────────────────
\newpage
\section*{Response to Peer Review}
% ──────────────────────────────────────────────────────────────
 
The table below summarises how the main remarks from reviews 96A, 96B, and 96C have
been addressed in this revised version.
 
\smallskip
\noindent
\footnotesize
\begin{tabular}{p{0.22\textwidth}p{0.18\textwidth}p{0.52\textwidth}}
\toprule
\textbf{Remark} & \textbf{Source} & \textbf{Action taken} \\
\midrule
Validation of inclusion/exclusion criteria &
96A, 96B &
Added paragraph in \S5.1.1: two team members independently inspect 30 sampled repositories after automated filtering; disagreements lead to exclusion. \\
\addlinespace
Data preprocessing steps &
96B &
Added paragraph in \S5.1.2 detailing three steps: exclusion of unparseable YAML files, deduplication across time points, and UTF-8/Unix normalisation. \\
\addlinespace
Justification of sample size &
96A, 96C &
Added justification in \S5.1.3 comparing our 60-file sample to prior studies~\cite{gallaba2020,vassallo2020b}; limited power acknowledged as a boundary condition. \\
\addlinespace
Confirmation bias mitigation &
96C &
Strengthened \S5.1.3 with two additions: (1) structured binary checklist to limit interpretive freedom; (2) audit of 10 files by a fourth reviewer external to the study. \\
\addlinespace
CI-Odor comparability &
96C &
Clarified in \S5.3 that CI-Odor results are reported in a dedicated subsection and not merged into the precision/recall/F1 table used for Gretel and CD-Linter. \\
\addlinespace
Statistical test assumptions &
96B &
Added sentence in \S5.2 stating that both Mann--Whitney U assumptions (ordinal data; platform independence) will be verified before application. \\
\addlinespace
Summary table of related work &
96A, 96C &
Added Table~1 in \S3.1 comparing prior studies on platform coverage, anti-patterns, tool evaluated, and cross-platform scope. \\
\addlinespace
Repetition of novelty claim &
96C &
Trimmed the \emph{State-of-the-art positioning} paragraph in \S1; redundant restatements replaced by a forward reference to \S3. \\
\bottomrule
\end{tabular}
 
\normalsize
\smallskip
\paragraph{Remarks not acted upon.}
\emph{Adding a pipeline diagram (96C):} the structured subsections of \S5 already
describe the protocol sequentially; a figure would duplicate rather than clarify.
\emph{Course structure compliance (96A):} the current structure follows standard
empirical software engineering registered report conventions, which we consider more
appropriate for the scientific content; all required course elements remain present.
 
\end{document}
